% ENTER27 长摘要。正文只收录 Gemini 对话中的英文定稿，不含中文修改说明。
% Springer LNCS（llncs）。双盲审：不恢复作者姓名与单位。
% 会议报告类扩展摘要通常限 2 页（含参考文献）。本文件按英文定稿排版，不删节。
\documentclass[runningheads]{llncs}
\usepackage[T1]{fontenc}
\usepackage{newtxtext}
\usepackage[varvw]{newtxmath}

\begin{document}

\title{Who Represents the Destination, and What Do They Evaluate? Social-Media Posts about Qingdao}
\titlerunning{Social-Media Posts about Qingdao}

\author{Anonymous Author}
\authorrunning{Anonymous Author}
\institute{Author names and affiliations are omitted for double-blind review.}

\maketitle

\begin{abstract}
Social-media posts are often read as direct proxies for tourists' destination image. Treating posts as produced discourse rather than static cognitive images, we examine how institutional accounts and individual creators portray a single destination, Qingdao. After cleaning, 10{,}178 of 11{,}599 Qingdao tourism-related posts from four platforms were retained as destination-relevant. In a pilot of 100 posts (476 segments), practical trip-planning information (58\% of segments) outweighed affective and sensory expression (30\%). Core attractions were evaluated overwhelmingly favourably (155 of 166 evaluations), whereas operational aspects (price, service, accessibility, and crowding) were evenly split (20 favourable vs.\ 19 unfavourable).

\keywords{destination image \and user-generated content \and content creators \and content analysis.}
\end{abstract}

\section{Background}
Destination image is an interactive system of thoughts, opinions, feelings, visualizations, and intentions toward a destination~\cite{tasci2007}, which content frequency counts capture only in part. Image forms through induced, autonomous, and organic information agents~\cite{gartner1994}; social-media posts now act as such agents and can be coded for cognitive and affective attributes~\cite{garay2019}, but an online post is not equivalent to the internal image a visitor holds. Studies such as~\cite{sun2021} compare official marketing (projected image) with tourist-generated content (perceived image) by the themes each contains. Less is known about the functions a post performs, which attributes it evaluates, and how creator identity shapes both.

\section{Purpose}
We compare how institutional accounts and individual creators represent Qingdao: their communicative strategies, the resources and attributes they evaluate, and how these co-occur within text segments. We hypothesize that, even when depicting identical destination resources, institutional posts stay uncritically promotional, whereas individual creators balance positive portrayals with practical caveats, critiques, or advice.

\section{Method}
We employed theory-driven structured content analysis~\cite{krippendorff2019}. Before coding, embedding-based semantic segmentation (Qwen3-Embedding-8B) under fixed boundary rules divided each post into self-contained units of meaning; coders did not alter the boundaries. Each segment was coded for communicative strategy, target destination resource, evaluative attribute, and evaluative valence and intensity. Account metadata were coded separately, and creator grouping criteria were fixed before any group comparison. A relevance classifier trained on 1{,}300 manually labelled posts screened the 11{,}599 Qingdao tourism-related posts collected from four platforms and retained 10{,}178 as destination-relevant. Two trained coders independently coded the pilot (Cohen's $\kappa = 0.82$ for evaluative direction) and resolved disagreements by discussion.

\section{Results to Date}
\paragraph{What posts do.}
The pilot covers 100 randomly sampled posts (476 segments). Practical information (itineraries, pricing, how-to advice) appears in 58\% of segments, sensory or emotional expression in 30\%, background knowledge in 11\%, and calls for interaction in 3\%. About half of the posts give explicit recommendations.

\paragraph{What posts evaluate.}
The coders agreed on both attribute and valence for 216 evaluations. Core attraction quality, activities, and destination atmosphere accounted for 166 evaluations (155 favourable vs.\ 10 unfavourable). By contrast, operational factors---price, service, facilities, accessibility, crowding, and management---accounted for 44, evenly split (20 favourable vs.\ 19 unfavourable).

\section{Significance}
Praise attaches mostly to endowed destination assets, whereas criticism concentrates on manageable operational touchpoints, so the actionable signal for destination marketing organizations (DMOs) lies in roughly one-fifth of evaluations, which aggregate sentiment scores obscure. In line with the conference theme (\emph{Shaping Tourism Technology: Hype, Human Centricity and ROI}), our workflow leaves high-volume filtering and segmentation to models and interpretive judgement to human coders. Social-listening tools that report sentiment by attribute and creator type would show DMOs where operational investment pays off.

\begin{thebibliography}{5}
\bibitem{garay2019}
Garay, L.: \#Visitspain. Breaking down affective and cognitive attributes in the social media construction of the tourist destination image. Tour. Manag. Perspect. \textbf{32}, 100560 (2019)

\bibitem{gartner1994}
Gartner, W.C.: Image formation process. J. Travel Tour. Mark. \textbf{2}(2--3), 191--216 (1994)

\bibitem{krippendorff2019}
Krippendorff, K.: Content Analysis: An Introduction to Its Methodology, 4th edn. SAGE, Los Angeles (2019)

\bibitem{sun2021}
Sun, W., Tang, S., Liu, F.: Examining perceived and projected destination image: A social media content analysis. Sustainability \textbf{13}(6), 3354 (2021)

\bibitem{tasci2007}
Tasci, A.D.A., Gartner, W.C., Cavusgil, S.T.: Conceptualization and operationalization of destination image. J. Hosp. Tour. Res. \textbf{31}(2), 194--223 (2007)
\end{thebibliography}

\end{document}
